% Complete final Kamp argument with explicit local scopes and finite windows.

\subsection{Present progressive: definitions and complete nonexpressibility proof}
\label{sec:annex-p09-kamp-present-progressive}

\begin{annexitem}{annex:p9.kamp.progressive.definitions}{thm:p09-kamp-progressive-definitions}{The two definitions and their exact boundary behavior}
\begin{proof}
If \(\rho(t_0)\) holds at \(x\), its two witnesses
\(a<x<b\), its current-point conjunct and its two interval
universals put every point of \([a,b]\) in the input.
Indeed each point is an endpoint, lies strictly between
\(a,x\), is \(x\), or lies strictly between \(x,b\).
Conversely, a closed interval \([a,b]\subseteq A\) with
\(a<x<b\) supplies every conjunct of \(\rho\). This
proves the displayed set characterization of \(\mathsf{PR}\),
which also agrees with the operation announced on p.~19.

Forgetting the witness-predicate conjuncts in \(\rho\)
gives \(\rho_0\), and both formulas require the current
point to be in \(A\). Both require a strictly earlier
and a strictly later witness. Hence both are false at a
first point, a last point, or the sole point of a
one-point order.

On a dense order, suppose \(\rho_0\) holds at \(x\)
with witnesses \(a<x<b\). Choose \(a'\) strictly
between \(a,x\) and \(b'\) strictly between \(x,b\).
The two universal clauses of \(\rho_0\) make
\(a',b'\) members of \(A\), and give the required
input values throughout their open intervals to \(x\).
Thus \(a',b'\) witness \(\rho\). The other implication
was already proved. This applies whenever \(\rho_0\)
has witnesses, and both formulas are false at any endpoints,
so equality holds on the whole dense order.

Now take a discrete order with the source's conditional
immediate-neighbor property. If \(x\) has points on
both sides, let \(x^-\) and \(x^+\) be its immediate
neighbors. The open intervals \((x^-,x)\) and
\((x,x^+)\) are empty. Thus \(x\in A\) alone
suffices for \(\rho_0\), with those two witnesses.
Its converse follows from the explicit current-point
conjunct. When a neighbor is absent, the corresponding
strict existential is false. This proves
\(\mathsf{PR}_0(A)=A\cap T_*\).
The formula \(\neg(q_0\land\neg q_0)\) is true at
every time point. Its past and future assert the existence
of an earlier and a later point, respectively. Their
conjunction with \(q_0\) therefore expresses precisely
\(A\cap T_*\), including at endpoints; when the order
has no endpoints these two side conditions are always true.

For \(\rho\), any past witness \(a<x\) either is
\(x^-\) or lies below it. In the first case its witness
conjunct gives \(x^-\in A\); in the second case the
past interval universal gives that membership. The same
reasoning on the right gives \(x^+\in A\). Conversely,
if \(x^-,x,x^+\in A\), the two immediate neighbors
are permitted witnesses and both interval universals
are vacuous. This proves the three-point characterization.
On \(\mathbb Z\), a singleton input consequently has
the stated weak and revised values. At the first point
of \(\mathbb N\), neither formula has a past witness,
including when that point belongs to the input.

Finally, by the source's strict argument order,
\(S(q_0,q_0)\) expresses the past existential and
universal of \(\rho\), and \(U(q_0,q_0)\) expresses
its future existential and universal. Conjoining the
current atom gives exactly \(\rho\), with no order
assumption beyond linearity. This proves the stated
\(\mathsf{TL}_1\) representation on all orders.
The linguistic interpretation motivating either formula
is separate from these mathematical equivalences
\cite[pp.~19, 112--113]{Kamp1968}.
\end{proof}
\end{annexitem}

\begin{annexitem}{annex:p9.kamp.progressive.two.sided}{thm:p09-kamp-progressive-two-sided}{The two-sided construction without a global cofinality assumption}
\begin{proof}
The set \(I\) is convex, and its markers are cofinal
in both directions within \(I\) by definition. Outside
the filled interval, the even markers belong to \(A\)
and the odd markers do not. For indices arbitrarily
far in either direction the markers are outside that
fixed interval. Thus both \(A\cap I\) and
\(I\setminus A\) have points strictly on each side
of every point of \(I\).

We prove the four-choice invariant by structural
induction. The atomic truth set restricts to
\(A\cap I\). Complement and intersection preserve
the field \(\{\varnothing,I,A\cap I,I\setminus A\}\):
complement swaps its two complementary pairs, and
intersection either retains a selected piece or removes
it. Suppose the induction hypothesis holds for an
operand \(\xi\). If its restriction to \(I\)
is any of the three nonempty choices, that restriction
is cofinal in both directions in \(I\). Hence every
current point of \(I\) has an earlier and a later
\(\xi\)-witness inside \(I\). Both \(P\xi\)
and \(F\xi\) are true throughout \(I\).

If instead \(\xi\) is false throughout \(I\),
a past witness for a point of \(I\) must lie
outside \(I\) on its lower side, and a future
witness must lie on its upper side. Convexity makes
every exterior point lower than all of \(I\) or
higher than all of \(I\). Existence of a witness
on each side is therefore independent of the current
point in \(I\). Each modality is either true
throughout \(I\) or false throughout it. This
also yields a permitted field member. The two
modal cases together with the Boolean cases exhaust
the \(\mathsf{TL}_0\) syntax.

We next compute the exact progressive image. Every
point strictly between \(a_{-2},a_2\) has these
two input points as witnesses, with their whole
closed interval in \(A\). Thus the open interval
is included in \(\mathsf{PR}(A)\). Its left
endpoint cannot have a valid earlier witness:
the odd marker \(a_{-3}\) is false, all earlier
input points are at most \(a_{-4}\), and no
input point lies strictly between \(a_{-3}\)
and \(a_{-2}\). Its right endpoint fails the
future requirement by the corresponding false
marker \(a_3\).

Every input point outside the filled interval is
an even marker. Immediately before it in the
marker ordering is a false odd marker, with no
input point between that odd marker and the
given even point. This prevents an all-input
closed interval from an earlier input point to
the given point. For a marker on the left side
of the filled block the same fact holds, and
for a marker on the right side the intervening
odd marker is outside the block as well. Thus
none of these points is progressive. A point
outside \(A\) already fails the current-atom
conjunct. This proves
\(\mathsf{PR}(A)=(a_{-2},a_2)\).

Both \(a_0,a_4\) belong to \(A\cap I\),
but only \(a_0\) belongs to this image. Every
set in the four-choice invariant has the same
membership at those two points, so no candidate
formula defines the progressive operation on
this interpretation. Only cofinality within
\(I\) was used; values of operand formulas
at other time points were explicitly allowed
in the modal induction
\cite[pp.~113--114]{Kamp1968}.
\end{proof}
\end{annexitem}

\begin{annexitem}{annex:p9.kamp.progressive.one.sided}{thm:p09-kamp-progressive-one-sided}{Both one-sided inductions and the exact block images}
\begin{proof}
First use the increasing sequence. The set \(I_+\)
is convex and begins at \(a_0\). Each of its
points is below some later marker, so the
markers are cofinal in \(I_+\), whether
or not they are cofinal in \(T\). For a fixed
even \(i\), the input \(A_i^+\) contains
every sufficiently late even marker and excludes
every sufficiently late odd marker. Both colors
are therefore cofinal in every local tail
\(H_m^+\).

The definition of depth shows by induction that
every value of \(d_-\) is even: its base is
zero, Boolean cases retain or maximize even
values, future leaves the value unchanged,
and past adds two. We now prove the claimed
truth-set invariant by induction on \(\eta\),
for every even \(i\geq d_-(\eta)\).

For the atom the restriction is exactly
\(H_0^+\cap A_i^+\). Negation takes complement
inside the same local tail. For a conjunction,
restrict both operand invariants to the common
smaller tail at the maximum of their depths.
Each restriction is still empty, full, input,
or complemented input on that tail. Intersecting
the two retains one of these four possibilities:
empty absorbs, full leaves the other unchanged,
equal colors intersect to themselves, and
opposite colors have empty intersection.

Suppose \(\eta=F\xi\), and put
\(m=d_-(\xi)=d_-(\eta)\). The induction
hypothesis applies on \(H_m^+\). If the
operand restriction is nonempty, it is either
the full tail or one of its two cofinal colors.
Every point of the tail therefore has a later
operand witness within the tail. The output
is the full tail. If the operand is empty
there, a later witness for a point of
\(H_m^+\) cannot lie in the earlier part
of \(I_+\). It must lie above all of
\(I_+\), outside that convex set. Existence
of such a witness is independent of the
current point of the local tail. Thus this
case also gives either the empty or full tail.
This is why no increment is needed for future
in this directional depth.

Suppose \(\eta=P\xi\), put
\(m=d_-(\xi)\), and recall
\(d_-(\eta)=m+2\). The assumption gives
\(i\geq m+2\). The two markers
\(a_{m+1},a_{m+2}\) lie in the old tail
\(H_m^+\). The first is odd and outside
the filled interval, so it is not in the
input; the second is even and is in the
input. The filled interval starts no earlier
than \(a_{m+4}\), so these colors are
unaffected by the filling. Both markers are
at or below \(a_{m+2}\).

If the operand restriction on \(H_m^+\)
is full or is one of the two colors, at
least one of these two markers satisfies
the operand. It is strictly before every
current point in \(H_{m+2}^+\), so
\(P\xi\) is true on the whole new tail.
The witness at \(a_{m+2}\) is allowed:
the current point is strictly later than
that threshold. If the operand is empty
on the old tail, every possible past
witness for a point of the new tail is
at or below \(a_m\), including points
outside \(I_+\) on the left. Existence
of any such witness again gives a constant
truth value on the new tail. Points outside
\(I_+\) on the right cannot be past
witnesses. This proves the past step,
including operand values at all time points,
not just the selected markers.

All primitive cases are proved. To compute
the progressive image, the two endpoints
of the filled interval witness every point
of \((a_{i+2},a_{i+4})\). The endpoints
have even indices. The false odd marker
\(a_{i+1}\) prevents the left endpoint
from having an all-input past interval;
\(a_{i+5}\) prevents the right endpoint
from having an all-input future interval.
Any input point outside the block is an
even marker separated from all earlier
input points by the preceding false odd
marker, except \(a_0\), which has no
earlier input point at all. Hence none
of these points is progressive. This
gives exactly the stated open interval.

For any candidate formula, choose even
\(i\geq d_-(\eta)\). The points
\(a_{i+3}\) and \(a_{i+6}\) both
belong to the input and to the relevant
tail. The former is inside the filled
interval, while the latter is an exterior
even marker. Their progressive values
differ, whereas each of the four possible
formula truth sets agrees at them. Thus
this interpretation prevents a definition.

We give the full order-reversal justification
for the decreasing case. Let
\(\mathcal J^{\mathrm{op}}=\langle T,>\rangle\),
and let \(\eta^*\) be obtained by interchanging
\(P,F\) everywhere, retaining the atom and
Boolean constructors. For every input, truth
of \(\eta\) in \(\mathcal J\) equals
truth of \(\eta^*\) in \(\mathcal J^{\mathrm{op}}\)
at every point. The atomic case retains the
same input set. Negation and conjunction
preserve equal truth values. An earlier
witness in \(\mathcal J\) is precisely
a later witness in its reverse, giving
the case for \(P\); a later witness is
an earlier witness in the reverse, giving
the case for \(F\). These are all cases
of the syntax induction. The depth recursion
then gives \(d_-(\eta^*)=d_+(\eta)\).

The sequence \((b_n)\) is increasing
in the reverse order. Its increasing
local domain and tails are exactly
\(I_-\) and \(H_m^-\) in the original
order, and its closed filled interval is
\([b_{i+4},b_{i+2}]\). The progressive
operation is itself invariant under order
reversal, since a witnessing closed interval
has the same set of points with its
endpoints exchanged. Applying the already
proved increasing result therefore gives
the stated decreasing invariant, exact
image, and nonexpressibility comparison,
with future adding two to \(d_+\) and
past adding zero. This proves both one-sided
source constructions in their full local
scopes
\cite[pp.~114--116]{Kamp1968}.
\end{proof}
\end{annexitem}

\begin{annexitem}{annex:p9.kamp.progressive.window}{thm:p09-kamp-progressive-window}{Every case of the two-depth invariant on finite or infinite chains}
\begin{proof}
Fix the indicated chains and even \(i,j\).
The two filled intervals are disjoint, since
every left marker is below every right marker.
We prove the invariant by structural induction,
for all formulas whose two depths are at most
\(i,j\), respectively. Every marker needed
in the proof exists in both the finite and
infinite versions of the construction.

For \(q_0\) the restricted truth set is
the input intersected with \(C_{00}\).
Negation takes complement in the same
window. For a conjunction, the two
maximum depths give a window contained
in each operand's window: increasing
the left index moves the lower bound
right, and increasing the right index
moves the upper bound left. Restriction
of an operand's four choices to this
common window leaves the same four
types. Their intersections are again
one of the four, by the complete
Boolean cases already described.

Suppose \(\eta=P\xi\). Put
\(m=d_-(\xi),n=d_+(\xi)\).
The old window is \(C_{mn}\), and
the new one is \(C_{m+2,n}\).
Since \(i\geq m+2\), the markers
\(a_{m+1},a_{m+2}\) are strictly
before the left filled interval. Their
indices are odd and even respectively,
so the first is outside \(A_{ij}\)
and the second belongs to it. Neither
can belong to the right filled interval,
which is above every left marker.
They lie in the old window and at or
below the new window's lower boundary.

If the operand is the full old window
or either input color there, one of
these two points satisfies it. That
point is strictly earlier than every
point of the new window. Hence
\(P\xi\) is true throughout the
new window. If the operand is empty
in the old window, any earlier
operand witness for a new-window
point must be at or below \(a_m\).
Conversely any such witness is before
every point of the new window.
Existence of one is therefore a
constant condition there. This case
gives either the empty or full new
window. All possible values outside
the old window have been included
in this argument.

For \(\eta=F\xi\), retain the same
old indices \(m,n\). The new window
is \(C_{m,n+2}\), and
\(j\geq n+2\). The right markers
\(b_{n+1},b_{n+2}\) are outside
the right filled interval, with the
first outside the input and the
second inside. They cannot be in
the left filled interval. They
lie in the old window and at or
above the new upper boundary.
For a full or colored operand
restriction, one of them is a
strictly later witness for every
new-window point. For an empty
operand restriction, any future
witness must be at or above
\(b_n\), a condition independent
of the current point in the new
window. Thus both future subcases
also give either empty or full
output. This proves the full
modal induction, without assuming
anything about unselected points
outside the central window.

We compute the exact progressive image.
Every point in either displayed open
filled interval has its two filled
endpoints as witnesses. At the left
filled interval's lower endpoint,
the false marker \(a_{i+1}\) blocks
the past requirement. At its upper
endpoint, \(a_{i+5}\) blocks the
future requirement. The right block's
upper endpoint is blocked on the
future side by \(b_{j+1}\), and
its lower endpoint is blocked on
the past side by \(b_{j+5}\).
All these marker indices are odd,
and all the markers are outside
both filled intervals.

Any remaining input point is an
even marker outside the two blocks.
For a left-chain marker of positive
even index, the preceding odd marker
separates it from all earlier input
points, with no input point between
that odd marker and the given point.
The first left marker has no earlier
input point. Hence none of these
left-chain points is progressive.
For a right-chain marker, use the
corresponding odd marker on its
later side; the first right marker
has no later input point. This
also treats the innermost markers
of a finite pair of chains: the
left one fails its past condition
and the right one its future
condition, even if those two
markers happen to be adjacent
in the whole order. Noninput
points fail the present conjunct.
No other progressive point exists.

Since \(i\geq d_-(\eta)\), the
left filled interval is strictly
above \(a_{d_-(\eta)}\). It
is below every right marker,
and hence below \(b_{d_+(\eta)}\).
The right filled interval is
inside the same window by
\(j\geq d_+(\eta)\).
The two points \(a_{i+3},a_{i+6}\)
also lie in that window. The first
is in the filled input interval
and in its progressive image;
the second is an even input
marker strictly outside the
block and has no progressive
membership. This completes
both the invariant and every
asserted image statement for
finite and infinite chains
\cite[p.~117]{Kamp1968}.
\end{proof}
\end{annexitem}

\begin{annexitem}{annex:p9.kamp.present.progressive}{thm:p09-kamp-present-progressive}{The unrestricted infinite-order conclusion and boundary examples}
\begin{proof}
Let \(\eta\) be any proposed one-input
\(\mathsf{TL}_0\)-definition on the given
infinite strict linear order. Put
\[
 i=d_-(\eta),\qquad j=d_+(\eta),\qquad
 N=\max(i,j)+6.
\]
The two depths are even nonnegative
integers. Choose \(2(N+1)\) distinct
points of \(T\), which is possible
because \(T\) is infinite, and sort
this finite set in its linear order.
Label its first \(N+1\) members in
increasing order as \(a_0,\ldots,a_N\).
Label the remaining members in
decreasing order as \(b_0,\ldots,b_N\).
They satisfy every separation and
index condition of the finite-window
construction. This selection is
finite; no cofinal or infinite
monotone sequence is required.

Interpret \(q_0\) by \(A_{ij}\).
The window lemma places the truth
set of \(\eta\), restricted to
\(C_{ij}\), in its four-choice
field. Each member of that field
has equal membership at
\(a_{i+3},a_{i+6}\), because
both points belong to the input
and the window. The exact
progressive image contains the
first and excludes the second.
Thus \(\eta\) fails to define
\(\mathsf{PR}\) on this input.
Since the proposed formula was
arbitrary, no finite formula
in the stated language defines
the operation on all inputs.

This selection is available in
an infinite order with neither
endpoint, with only a first
endpoint, with only a last
endpoint, or with both. The
first-order definition of
\(\mathsf{PR}\) continues to
make it false at any actual
endpoint. The comparison points
are strictly internal selected
markers with other markers on
both sides, so that endpoint
qualification is preserved in
every case. Finite truth-functional
abbreviations do not enlarge
the basis, because they expand
into negation and conjunction
using the same present atom.

For the explicit \(\mathbb N\)
example with \(i=0\), the
input contains the consecutive
points two, three, four and
otherwise just the even points.
The discrete three-point
criterion makes precisely
three progressive. If instead
\(i=1\), the filled interval
adds three and five to the
even points, producing the
five-point consecutive block
two through six. Its interior
three, four, five has both
neighbors in the input; every
other input point is missing
at least one input neighbor
or lies at the first time
point. This gives exactly
\(\{3,4,5\}\), verifying
why the even-index restriction
is needed for the simpler
exact-image formula.

The order \(\omega+1\) has
its least natural number as
first point and its added
limit point as last point.
It has no infinite strictly
decreasing sequence: after
at most the first selected
point such a sequence would
lie among natural numbers,
and a natural number has
only finitely many smaller
natural numbers. Thus one
cannot infer two separated
infinite chains merely from
the existence of both
endpoints. The finite
construction above still
works in this order, using
any sufficiently long finite
list of natural numbers
below its last point. This
verifies the bounded infinite
case without the unsupported
sequence premise and completes
the source's announced
infinite-order nonexpressibility
claim
\cite[pp.~19, 112--117]{Kamp1968}.
\end{proof}
\end{annexitem}
